\section{Introduzione}
\label{sec:introduzione}

\subsection{Il non-equilibrio termochimico nel rientro atmosferico}
\label{sec:contesto}

Un veicolo che rientra nell'atmosfera a velocità ipersonica genera un urto
forte, dietro il quale la temperatura sale a decine di migliaia di kelvin: i
modi vibrazionali si eccitano, i livelli elettronici si popolano e le molecole
si dissociano. Questi processi non raggiungono l'equilibrio alla stessa
velocità: traslazione e rotazione si equilibrano in pochi urti molecolari, la
vibrazione in molti di più, e il gas dietro l'urto si trova in non-equilibrio
termico. Il modello a due temperature di Park~\cite{park1990,park1993} lo
descrive tenendo distinte una temperatura traslazionale-rotazionale e una
vibro-elettronica, legate da un tempo di rilassamento. A questo si somma
l'accoppiamento con la chimica: una molecola eccitata vibrazionalmente si
dissocia più facilmente, e le reazioni a loro volta sottraggono o restituiscono
energia vibrazionale. Dalla descrizione di questi processi dipende la
previsione dei carichi termici sul veicolo.

\subsection{Il modello a una temperatura e i suoi limiti}
\label{sec:una-temperatura}

Una molecola può immagazzinare energia in modi diversi: spostandosi nello
spazio (traslazione), ruotando su se stessa (rotazione), facendo oscillare i
propri atomi (vibrazione) ed eccitando i propri elettroni. Nella maggior parte
delle simulazioni CFD, compresi i solver standard di OpenFOAM, il gas è
descritto da una sola temperatura: si assume che l'energia ricevuta dal gas si
ripartisca subito fra tutti questi modi nelle proporzioni di equilibrio, per
cui un solo numero basta a sapere quanta energia c'è in ciascuno, e con lo
stesso numero si calcola anche la velocità delle reazioni chimiche. È
un'ipotesi eccellente nella gran parte delle applicazioni, perché le molecole
si urtano così spesso che la ripartizione avviene molto più in fretta di
quanto cambi il flusso.

Dietro un urto ipersonico questa ipotesi cade: l'energia arriva prima al moto
delle molecole e solo dopo molti urti alla vibrazione, in un tempo paragonabile
a quello che il gas impiega ad attraversare la regione dietro l'urto. Un
modello a una temperatura non vede questo ritardo: prevede appena dietro l'urto
una temperatura più bassa di quella reale e salta la zona in cui il gas si
porta gradualmente all'equilibrio. Lo stesso accade per la chimica: la
dissociazione dipende sia dall'energia con cui le molecole si urtano sia da
quanto stanno già vibrando, due grandezze che una sola temperatura non può
distinguere. Il risultato sono errori sulla temperatura e sulla composizione
del gas, e quindi sulla posizione dell'urto e sul calore che raggiunge la
parete~\cite{park1990}.

\subsection{La soluzione: il modello a due temperature}
\label{sec:soluzione}

Il modello a due temperature di Park~\cite{park1990,park1993} supera questi
limiti usando due temperature invece di una: la prima descrive il moto e la
rotazione delle molecole, che si equilibrano quasi subito, la seconda la
vibrazione e gli elettroni, che restano indietro. Subito dietro l'urto le due
possono differire di migliaia di gradi; l'energia passa poi dall'una all'altra
gradualmente, tanto più in fretta quanto più il gas è caldo e denso, finché le
due temperature coincidono e il gas è tornato in equilibrio. Anche la chimica
le usa entrambe: la dissociazione è calcolata a una temperatura intermedia,
così che una molecola che vibra ancora poco si dissoci più lentamente, e
l'energia di vibrazione portata via dalle molecole che si spezzano viene tolta
da quella del gas. Resta un'approssimazione, perché dentro ciascuno dei due
gruppi l'energia è supposta in equilibrio, ma cattura l'essenziale senza dover
seguire ogni singolo livello energetico delle molecole.

\subsection{Il lavoro di riferimento}
\label{sec:riferimento}

Il lavoro di riferimento è l'articolo di Casseau
\textit{et al.}~\cite{casseau2016}, dell'Università di Strathclyde, che
presenta \hyfoam, un solver a due temperature per flussi ipersonici reagenti
costruito su OpenFOAM e distribuito come codice aperto~\cite{hystrath}. La
prima parte dell'articolo verifica il solver su una serie di bagni termici:
piccoli volumi di gas fermi e isolati, preparati fuori equilibrio e lasciati
evolvere, in cui ogni processo si può studiare da solo. I casi vanno dall'azoto
puro, con e senza energia elettronica, alle miscele di azoto con azoto atomico
e con ossigeno, fino ai gas che reagiscono e all'aria a cinque specie, e i
risultati sono confrontati con altri codici CFD e con codici DSMC, che
simulano il gas seguendo il moto e gli urti di un gran numero di particelle.
Sono questi i casi che il presente lavoro riproduce.
